% Options for packages loaded elsewhere
% Options for packages loaded elsewhere
\PassOptionsToPackage{unicode}{hyperref}
\PassOptionsToPackage{hyphens}{url}
\PassOptionsToPackage{dvipsnames,svgnames,x11names}{xcolor}
%
\documentclass[
  letterpaper,
  DIV=11,
  numbers=noendperiod]{scrartcl}
\usepackage{xcolor}
\usepackage{amsmath,amssymb}
\setcounter{secnumdepth}{-\maxdimen} % remove section numbering
\usepackage{iftex}
\ifPDFTeX
  \usepackage[T1]{fontenc}
  \usepackage[utf8]{inputenc}
  \usepackage{textcomp} % provide euro and other symbols
\else % if luatex or xetex
  \usepackage{unicode-math} % this also loads fontspec
  \defaultfontfeatures{Scale=MatchLowercase}
  \defaultfontfeatures[\rmfamily]{Ligatures=TeX,Scale=1}
\fi
\usepackage{lmodern}
\ifPDFTeX\else
  % xetex/luatex font selection
\fi
% Use upquote if available, for straight quotes in verbatim environments
\IfFileExists{upquote.sty}{\usepackage{upquote}}{}
\IfFileExists{microtype.sty}{% use microtype if available
  \usepackage[]{microtype}
  \UseMicrotypeSet[protrusion]{basicmath} % disable protrusion for tt fonts
}{}
\makeatletter
\@ifundefined{KOMAClassName}{% if non-KOMA class
  \IfFileExists{parskip.sty}{%
    \usepackage{parskip}
  }{% else
    \setlength{\parindent}{0pt}
    \setlength{\parskip}{6pt plus 2pt minus 1pt}}
}{% if KOMA class
  \KOMAoptions{parskip=half}}
\makeatother
% Make \paragraph and \subparagraph free-standing
\makeatletter
\ifx\paragraph\undefined\else
  \let\oldparagraph\paragraph
  \renewcommand{\paragraph}{
    \@ifstar
      \xxxParagraphStar
      \xxxParagraphNoStar
  }
  \newcommand{\xxxParagraphStar}[1]{\oldparagraph*{#1}\mbox{}}
  \newcommand{\xxxParagraphNoStar}[1]{\oldparagraph{#1}\mbox{}}
\fi
\ifx\subparagraph\undefined\else
  \let\oldsubparagraph\subparagraph
  \renewcommand{\subparagraph}{
    \@ifstar
      \xxxSubParagraphStar
      \xxxSubParagraphNoStar
  }
  \newcommand{\xxxSubParagraphStar}[1]{\oldsubparagraph*{#1}\mbox{}}
  \newcommand{\xxxSubParagraphNoStar}[1]{\oldsubparagraph{#1}\mbox{}}
\fi
\makeatother

\usepackage{color}
\usepackage{fancyvrb}
\newcommand{\VerbBar}{|}
\newcommand{\VERB}{\Verb[commandchars=\\\{\}]}
\DefineVerbatimEnvironment{Highlighting}{Verbatim}{commandchars=\\\{\}}
% Add ',fontsize=\small' for more characters per line
\usepackage{framed}
\definecolor{shadecolor}{RGB}{241,243,245}
\newenvironment{Shaded}{\begin{snugshade}}{\end{snugshade}}
\newcommand{\AlertTok}[1]{\textcolor[rgb]{0.68,0.00,0.00}{#1}}
\newcommand{\AnnotationTok}[1]{\textcolor[rgb]{0.37,0.37,0.37}{#1}}
\newcommand{\AttributeTok}[1]{\textcolor[rgb]{0.40,0.45,0.13}{#1}}
\newcommand{\BaseNTok}[1]{\textcolor[rgb]{0.68,0.00,0.00}{#1}}
\newcommand{\BuiltInTok}[1]{\textcolor[rgb]{0.00,0.23,0.31}{#1}}
\newcommand{\CharTok}[1]{\textcolor[rgb]{0.13,0.47,0.30}{#1}}
\newcommand{\CommentTok}[1]{\textcolor[rgb]{0.37,0.37,0.37}{#1}}
\newcommand{\CommentVarTok}[1]{\textcolor[rgb]{0.37,0.37,0.37}{\textit{#1}}}
\newcommand{\ConstantTok}[1]{\textcolor[rgb]{0.56,0.35,0.01}{#1}}
\newcommand{\ControlFlowTok}[1]{\textcolor[rgb]{0.00,0.23,0.31}{\textbf{#1}}}
\newcommand{\DataTypeTok}[1]{\textcolor[rgb]{0.68,0.00,0.00}{#1}}
\newcommand{\DecValTok}[1]{\textcolor[rgb]{0.68,0.00,0.00}{#1}}
\newcommand{\DocumentationTok}[1]{\textcolor[rgb]{0.37,0.37,0.37}{\textit{#1}}}
\newcommand{\ErrorTok}[1]{\textcolor[rgb]{0.68,0.00,0.00}{#1}}
\newcommand{\ExtensionTok}[1]{\textcolor[rgb]{0.00,0.23,0.31}{#1}}
\newcommand{\FloatTok}[1]{\textcolor[rgb]{0.68,0.00,0.00}{#1}}
\newcommand{\FunctionTok}[1]{\textcolor[rgb]{0.28,0.35,0.67}{#1}}
\newcommand{\ImportTok}[1]{\textcolor[rgb]{0.00,0.46,0.62}{#1}}
\newcommand{\InformationTok}[1]{\textcolor[rgb]{0.37,0.37,0.37}{#1}}
\newcommand{\KeywordTok}[1]{\textcolor[rgb]{0.00,0.23,0.31}{\textbf{#1}}}
\newcommand{\NormalTok}[1]{\textcolor[rgb]{0.00,0.23,0.31}{#1}}
\newcommand{\OperatorTok}[1]{\textcolor[rgb]{0.37,0.37,0.37}{#1}}
\newcommand{\OtherTok}[1]{\textcolor[rgb]{0.00,0.23,0.31}{#1}}
\newcommand{\PreprocessorTok}[1]{\textcolor[rgb]{0.68,0.00,0.00}{#1}}
\newcommand{\RegionMarkerTok}[1]{\textcolor[rgb]{0.00,0.23,0.31}{#1}}
\newcommand{\SpecialCharTok}[1]{\textcolor[rgb]{0.37,0.37,0.37}{#1}}
\newcommand{\SpecialStringTok}[1]{\textcolor[rgb]{0.13,0.47,0.30}{#1}}
\newcommand{\StringTok}[1]{\textcolor[rgb]{0.13,0.47,0.30}{#1}}
\newcommand{\VariableTok}[1]{\textcolor[rgb]{0.07,0.07,0.07}{#1}}
\newcommand{\VerbatimStringTok}[1]{\textcolor[rgb]{0.13,0.47,0.30}{#1}}
\newcommand{\WarningTok}[1]{\textcolor[rgb]{0.37,0.37,0.37}{\textit{#1}}}

\usepackage{longtable,booktabs,array}
\usepackage{calc} % for calculating minipage widths
% Correct order of tables after \paragraph or \subparagraph
\usepackage{etoolbox}
\makeatletter
\patchcmd\longtable{\par}{\if@noskipsec\mbox{}\fi\par}{}{}
\makeatother
% Allow footnotes in longtable head/foot
\IfFileExists{footnotehyper.sty}{\usepackage{footnotehyper}}{\usepackage{footnote}}
\makesavenoteenv{longtable}
\usepackage{graphicx}
\makeatletter
\newsavebox\pandoc@box
\newcommand*\pandocbounded[1]{% scales image to fit in text height/width
  \sbox\pandoc@box{#1}%
  \Gscale@div\@tempa{\textheight}{\dimexpr\ht\pandoc@box+\dp\pandoc@box\relax}%
  \Gscale@div\@tempb{\linewidth}{\wd\pandoc@box}%
  \ifdim\@tempb\p@<\@tempa\p@\let\@tempa\@tempb\fi% select the smaller of both
  \ifdim\@tempa\p@<\p@\scalebox{\@tempa}{\usebox\pandoc@box}%
  \else\usebox{\pandoc@box}%
  \fi%
}
% Set default figure placement to htbp
\def\fps@figure{htbp}
\makeatother





\setlength{\emergencystretch}{3em} % prevent overfull lines

\providecommand{\tightlist}{%
  \setlength{\itemsep}{0pt}\setlength{\parskip}{0pt}}



 


\KOMAoption{captions}{tableheading}
\makeatletter
\@ifpackageloaded{caption}{}{\usepackage{caption}}
\AtBeginDocument{%
\ifdefined\contentsname
  \renewcommand*\contentsname{Table of contents}
\else
  \newcommand\contentsname{Table of contents}
\fi
\ifdefined\listfigurename
  \renewcommand*\listfigurename{List of Figures}
\else
  \newcommand\listfigurename{List of Figures}
\fi
\ifdefined\listtablename
  \renewcommand*\listtablename{List of Tables}
\else
  \newcommand\listtablename{List of Tables}
\fi
\ifdefined\figurename
  \renewcommand*\figurename{Figure}
\else
  \newcommand\figurename{Figure}
\fi
\ifdefined\tablename
  \renewcommand*\tablename{Table}
\else
  \newcommand\tablename{Table}
\fi
}
\@ifpackageloaded{float}{}{\usepackage{float}}
\floatstyle{ruled}
\@ifundefined{c@chapter}{\newfloat{codelisting}{h}{lop}}{\newfloat{codelisting}{h}{lop}[chapter]}
\floatname{codelisting}{Listing}
\newcommand*\listoflistings{\listof{codelisting}{List of Listings}}
\makeatother
\makeatletter
\makeatother
\makeatletter
\@ifpackageloaded{caption}{}{\usepackage{caption}}
\@ifpackageloaded{subcaption}{}{\usepackage{subcaption}}
\makeatother
\usepackage{bookmark}
\IfFileExists{xurl.sty}{\usepackage{xurl}}{} % add URL line breaks if available
\urlstyle{same}
\hypersetup{
  pdftitle={Summative Assessment RM\&DA},
  pdfauthor={Your Name},
  colorlinks=true,
  linkcolor={blue},
  filecolor={Maroon},
  citecolor={Blue},
  urlcolor={Blue},
  pdfcreator={LaTeX via pandoc}}


\title{Summative Assessment RM\&DA}
\author{Your Name}
\date{2026-03-05}
\begin{document}
\maketitle

\renewcommand*\contentsname{Table of contents}
{
\hypersetup{linkcolor=}
\setcounter{tocdepth}{3}
\tableofcontents
}

\subsubsection{Short Brief}\label{short-brief}

\begin{itemize}
\item
  This is an open-book exam. You may use any resources available to you
  (notes, books, online documentation).
\item
  You are expected to work individually. Collaboration is not permitted.
\item
  All code and answers must be clearly presented in this document.
\item
  Ensure this document renders successfully to both an HTML and a PDF
  file.
\item
  Show all your work, including the R code used to arrive at your
  answers.
\end{itemize}

\section{Section 1: Data Wrangling}\label{section-1-data-wrangling}

This section uses the \texttt{dataset}, which contains data about
various car models.

\subsubsection{Step 1 - Load the
database}\label{step-1---load-the-database}

Load the dataset

\begin{Shaded}
\begin{Highlighting}[]
\CommentTok{\# Place your code here}
\end{Highlighting}
\end{Shaded}

Inspect the dataset for basic characteristics such as type of variables,
dimensions of the dataset, etc.

\begin{Shaded}
\begin{Highlighting}[]
\CommentTok{\# Place your code here}
\end{Highlighting}
\end{Shaded}

\subsubsection{Step 2 - Data wrangling}\label{step-2---data-wrangling}

\begin{itemize}
\tightlist
\item
  Create a new variable based on an existing one
\end{itemize}

The \texttt{am} column in the \texttt{mtcars} dataset represents the
transmission type (0 = automatic, 1 = manual). Create a new factor
variable called \texttt{transmission\_type} based on the am column to
acheive the result below:

\begin{verbatim}
                  am transmission_type
Mazda RX4          1            Manual
Mazda RX4 Wag      1            Manual
Datsun 710         1            Manual
Hornet 4 Drive     0         Automatic
Hornet Sportabout  0         Automatic
Valiant            0         Automatic
\end{verbatim}

\begin{itemize}
\tightlist
\item
  Create transmission\_type
\end{itemize}

\begin{Shaded}
\begin{Highlighting}[]
\CommentTok{\# Place your code here}
\end{Highlighting}
\end{Shaded}

\subsubsection{Step 3 - Data
Summarisation}\label{step-3---data-summarisation}

Calculate the mean mpg (miles per gallon) for cars with 4 cylinders
(cyl) and for cars with 6 cylinders.

\begin{itemize}
\tightlist
\item
  Calculate mean \texttt{mpg} for 4 and 6 cylinder cars
\end{itemize}

\begin{verbatim}
# A tibble: 2 x 2
    cyl mean_mpg
  <dbl>    <dbl>
1     4     26.7
2     6     19.7
\end{verbatim}

\begin{Shaded}
\begin{Highlighting}[]
\CommentTok{\# Place your code here}
\end{Highlighting}
\end{Shaded}

Question 4 ({[}Insert points here{]})

Filter the data to include only observations where the car has a manual
transmission and has more than 4 gears. How many observations meet these
criteria?

manual\_and\_more\_than\_4\_gears \textless- mtcars \%\textgreater\%
filter(am == 1 \& gear \textgreater{} 4)
nrow(manual\_and\_more\_than\_4\_gears)

Your answer here

Section 2: Data Visualization This section continues to use the mtcars
dataset.

Question 5 ({[}Insert Points Here{]})

Create a scatter plot showing the relationship between mpg and wt
(weight). Add appropriate labels for the axes and a title for the plot.

\section{Create scatter plot of mpg vs
wt}\label{create-scatter-plot-of-mpg-vs-wt}

ggplot(mtcars, aes(x = wt, y = mpg)) + geom\_point() + labs(title =
``Relationship between MPG and Weight'', x = ``Weight (1000 lbs)'', y =
``MPG'')

\emph{Your written answer for Question 5 goes here.}

Question 6 ({[}Insert Points Here{]})

Generate a boxplot to visualize the distribution of mpg across different
cyl (number of cylinders) values.

\section{Create boxplot of mpg by
cyl}\label{create-boxplot-of-mpg-by-cyl}

ggplot(mtcars, aes(x = factor(cyl), y = mpg)) + geom\_boxplot() +
labs(title = ``MPG Distribution by Number of Cylinders'', x = ``Number
of Cylinders'', y = ``MPG'')

\emph{Your written answer for Question 6 goes here.}

Question 7 ({[}Insert Points Here{]})

Create a bar chart showing the count of cars for each number of
cylinders (cyl).

\section{Create bar chart of cylinder
counts}\label{create-bar-chart-of-cylinder-counts}

ggplot(mtcars, aes(x = factor(cyl))) + geom\_bar() + labs(title = ``Car
Count by Number of Cylinders'', x = ``Number of Cylinders'', y =
``Count'')

\emph{Your written answer for Question 7 goes here.}

Section 3: Basic Statistics This section uses the mtcars dataset.

Question 8 ({[}Insert Points Here{]})

Perform an independent samples t-test to compare the mean mpg for cars
with automatic and manual transmissions. State the null and alternative
hypotheses, report the p-value, and interpret the results.

\section{Perform t-test}\label{perform-t-test}

t.test(mpg \textasciitilde{} am, data = mtcars)

\emph{Your written answer for Question 8 goes here.}

Question 9 ({[}Insert Points Here{]})

Calculate the Pearson correlation coefficient between mpg and hp
(horsepower). Interpret the strength and direction of the correlation.

\section{Calculate correlation}\label{calculate-correlation}

cor(mtcars\(mpg, mtcars\)hp)

\emph{Your written answer for Question 9 goes here.}

Question 10 ({[}Insert Points Here{]})

Conduct a chi-squared test to examine the relationship between cyl
(number of cylinders) and am (transmission type). State the null and
alternative hypotheses, report the chi-squared statistic and p-value,
and interpret the results.

\section{Perform chi-squared test}\label{perform-chi-squared-test}

table\_cyl\_am \textless- table(mtcars\(cyl, mtcars\)am)
chisq.test(table\_cyl\_am)

\emph{Your written answer for Question 10 goes here.}

Question 11 ({[}Insert Points Here{]})

Fit a linear model to predict mpg using wt (weight) and hp (horsepower).
Provide the R code, summarize the model, and interpret the coefficients
for wt and hp.

\section{Fit linear model}\label{fit-linear-model}

linear\_model\_mpg \textless- lm(mpg \textasciitilde{} wt + hp, data =
mtcars) summary(linear\_model\_mpg)

\emph{Your written answer for Question 11 goes here.}

Question 12 ({[}Insert Points Here{]})

Fit a logistic regression model to predict am (transmission type) using
mpg and wt (weight). Provide the R code, summarize the model, and
interpret the odds ratio for mpg.

\section{Fit logistic model}\label{fit-logistic-model}

logistic\_model\_am \textless- glm(am \textasciitilde{} mpg + wt, data =
mtcars, family = ``binomial'') summary(logistic\_model\_am)
exp(coef(logistic\_model\_am))

\emph{Your written answer for Question 12 goes here.}

Section 4: Basic Introduction to Multivariate Analysis This section uses
the iris dataset.

Question 13 ({[}Insert Points Here{]})

Briefly explain the purpose of Principal Component Analysis (PCA). In
the context of the `iris




\end{document}
